\scanpage{193}
\mspage{67}
% archive sheet number: 192

L'équivalence de a) et a$'$), \unsure{ou|resp.} de b) et b$'$), est $\pm$ tautologique. L'équivalence de a) et b) résulte aussitôt du NB inclus au \uline{1°}, compte tenu que les $\unsure{\pi|\Pi}$ tels que les objets de $\mathrm{Filt}^{\Pi}$ soient $\pi$-accessibles sont cofinaux dans l'ens des cardinaux $\in\unsure{\mathcal U|\mathscr U|U}$. (Cf 3.2, p. 16)

\margin{\rule{0.4pt}{3cm}}\uline{Proposition 5.13.} Toute réunion finie de parties $L$-$\pi$-accessibles de $\mathrm{Ob}\,\mathcal M$ est $L$-$\pi$-accessible. L'intersection d'une famille $(M_\alpha)_{\alpha\in A}$ de parties $L$-$\pi$-accessibles, avec $\mathrm{card}\,A\leq\pi$, est $L$-$\pi$-accessible.

\uline{Dém} La première assertion provient du fait qu'une réunion \uline{finie} de parties $L$-stables de $M$ est $L$-stable, \deleted{ce qui} et qu'on a \unsure{itou|idem} pour la notion de $\pi$-accessibilité. Les deux choses se prouvent d'ailleurs pratiquement de la même façon, en utilisant le

\uline{Lemme} 5.14. Soit $I$ un ens. ordonné filtrant, $(I_\alpha)_{\alpha\in A}$ une famille \uline{finie} \inserted{non vide} de parties de $I$, telle que $\bigcup I_\alpha=I$. Alors il existe $\alpha\in A$ tel que $I_\alpha$ soit cofinal, \deleted{dans $I$ et $=\ \alpha\in A$ et $i_0\in I$ tels que $I_\alpha\supset I/i_0$}

